\documentclass[12pt]{article}
\usepackage[utf8]{inputenc}
\usepackage[T1]{fontenc}
\usepackage{hyperref}
\title{Yeniche Language Resources for K-12 Curriculum Development}
\author{Compiled by Claude Sonnet 5.5 (claude-sonnet-5-5) for Andr\'as Kornai}
\date{}
\begin{document}
\maketitle

\section{Introduction}
Yeniche (Jenisch, Yenish; ISO 639-3: yec; Glottocode yeni1236) is a mixed language: a special vocabulary laid over German dialects. The vocabulary is mainly Rotwelsch, with Romani and Ashkenazic Hebrew elements. Sources describe a lexicon of a few hundred roots, used to mark group identity and keep private talk from outsiders. Each community has its own variety. Search results mention that the community has campaigned for recognition and that a Central Council of Jenisch People was founded in 2019. No official status in Germany was found.

\section{Primary Educational Resources}
No school materials were found in search results.

\section{Linguistic and Reference Materials}
\begin{itemize}
\item Engelbert Wittich, \emph{Die jenische Sprache}, a historical vocabulary on Project Gutenberg: \url{https://www.gutenberg.org/ebooks/50779}
\item Christian Efing, Wörterbuch of Lützenhardt Jenisch (PDF, Universität Wuppertal): \url{https://germanistik.uni-wuppertal.de/fileadmin/germanistik/Personal/Efing/Efing_Woerterbuch_LJ.pdf}
\item Overview page, Jenisch (Yenish), on the Kratylos site: \url{https://kratylos.org/~raphael/languageContact/ELA/languages/Jenisch.html}
\end{itemize}

\section{Religious and Community Resources}
None found in these results.

\section{Digital and Audio Resources}
No audio or software resources were found.

\section{Curriculum Development Context}
Yeniche is a lexical register rather than a full standalone grammar, and it is used mostly by adults. School instruction would therefore be in German, with Yeniche vocabulary taught as cultural and heritage content. The Wittich and Efing word lists could be seed material, but dialect variation and community privacy norms need to be respected. Community consultation would be essential.
\end{document}